\documentclass[../main.tex]{subfiles}
\ifSubfilesClassLoaded{\setcounter{chapter}{8}}{}
\begin{document}

\chapter{Modul, \texttt{pip}, Virtual Environment, dan Debugging}

\begin{subcpmk}
  \item Sub-CPMK 6.1: Mengorganisasi kode dengan \texttt{import}, mengelola paket dengan \texttt{pip}, dan menggunakan \texttt{venv} untuk lingkungan terisolasi
  \item Sub-CPMK 6.2: Melakukan \textit{debugging} sistematis dan menerapkan konvensi gaya penulisan kode yang konsisten
\end{subcpmk}

\SesiRPS{10}{6.1 dan 6.2}{$\sim$170 menit}{CPMK-6}

\section{Sarana dan prasarana}

Python 3.10+, terminal, akses \texttt{pip} sesuai kebijakan jaringan lab, editor/IDE.

\section{Prosedur kerja}

\begin{enumerate}[leftmargin=*]
  \item Ulangi contoh \texttt{import} dari pustaka standar.
  \item Buat \texttt{venv}, aktifkan, dan instal satu paket ringan (\texttt{pip install}).
  \item Latih pola \texttt{if \_\_name\_\_ == "\_\_main\_\_"} dan teknik debugging sederhana.
\end{enumerate}

\section{Mengimpor Modul}

\begin{lstlisting}[caption=Pola import (Tutorial Python ss6)]
import math                    # impor seluruh modul
from datetime import date      # impor spesifik
from os import path, getcwd    # impor beberapa nama
import random as rnd           # alias

print(math.sqrt(16))           # 4.0
print(date.today())
print(rnd.randint(1, 6))       # dadu acak
\end{lstlisting}

Hindari \texttt{from modul import *} pada kode produksi; gunakan impor eksplisit untuk menghindari konflik nama.

\section{Pola \texttt{if \_\_name\_\_ == "\_\_main\_\_"}}

Pola ini memungkinkan sebuah berkas Python dijalankan baik sebagai skrip utama maupun sebagai modul yang diimpor (Tutorial Python §6.1.1):

\begin{lstlisting}[caption=Pola if \textunderscore\textunderscore name\textunderscore\textunderscore == main]
# berkas: kalkulator.py

def tambah(a, b):
    return a + b

def kurang(a, b):
    return a - b

def main():
    a = float(input("a: "))
    b = float(input("b: "))
    print(f"Jumlah : {tambah(a, b)}")
    print(f"Selisih: {kurang(a, b)}")

if __name__ == "__main__":
    # Blok ini hanya dijalankan jika file dieksekusi langsung,
    # BUKAN saat diimpor oleh modul lain
    main()
\end{lstlisting}

\begin{lstlisting}[caption=Mengimpor fungsi dari berkas lain menggunakan modul]
# berkas: tes.py
from kalkulator import tambah, kurang
# Blok if __name__ == "__main__" di kalkulator.py TIDAK berjalan

print(tambah(3, 4))   # 7
\end{lstlisting}

\section{Menjelajahi Modul dengan \texttt{dir()}}

Fungsi \texttt{dir()} menampilkan daftar nama yang didefinisikan dalam sebuah modul atau objek (Tutorial Python §6.3):

\begin{lstlisting}[caption=Menjelajahi modul dengan dir()]
import math

# Lihat semua nama dalam modul math
print(dir(math))
# ['__doc__', '__name__', ..., 'cos', 'exp', 'floor', 'log', ...]

# Cari fungsi yang mengandung 'log'
fungsi_log = [nama for nama in dir(math) if "log" in nama]
print(fungsi_log)   # ['log', 'log10', 'log1p', 'log2']

# Tanpa argumen: daftar nama di scope saat ini
x = 10
print(dir())   # akan mencantumkan 'math', 'x', dll.
\end{lstlisting}

\section{Pustaka Standar yang Berguna}

Python hadir dengan pustaka standar yang kaya. Berikut beberapa modul penting (Tutorial Python §10):

\begin{lstlisting}[caption=Contoh modul pustaka standar]
import os
import pathlib
import random
import math
import json
import datetime
import sys

# os dan pathlib -- berkas dan direktori
print(os.getcwd())                          # direktori saat ini
print(pathlib.Path.home())                  # folder home user
print(os.path.exists("data.txt"))           # cek keberadaan berkas

# random -- pembangkit acak
print(random.randint(1, 100))               # integer acak
print(random.choice(["apel", "jeruk"]))     # pilihan acak
random.shuffle([1, 2, 3, 4, 5])            # acak urutan in-place

# math -- fungsi matematika
print(math.pi, math.e)
print(math.factorial(10))
print(math.gcd(48, 18))                     # FPB

# datetime -- tanggal dan waktu
sekarang = datetime.datetime.now()
print(sekarang.strftime("%d/%m/%Y %H:%M"))

# sys -- informasi interpreter
print(sys.version)
print(sys.platform)
\end{lstlisting}

\section{Virtual Environment dengan \texttt{venv}}

\textit{Virtual environment} adalah direktori terisolasi yang berisi instalasi Python sendiri, sehingga dependensi antar proyek tidak saling mengganggu (Tutorial Python §12):

\begin{lstlisting}[language=bash,caption=Membuat dan mengaktifkan virtual environment]
# Membuat virtual environment di folder .venv
python -m venv .venv

# Mengaktifkan (Windows)
.venv\Scripts\activate

# Mengaktifkan (Linux/macOS)
source .venv/bin/activate

# Menonaktifkan
deactivate
\end{lstlisting}

\begin{catatan}
Setelah virtual environment aktif, prompt shell berubah menampilkan nama venv (misalnya \texttt{(.venv)}). Semua instalasi \texttt{pip} dalam kondisi aktif ini hanya mempengaruhi venv, bukan instalasi Python global.
\end{catatan}

\section{Mengelola Paket dengan \texttt{pip}}

\texttt{pip} adalah manajer paket Python (Tutorial Python §12.3):

\begin{lstlisting}[language=bash,caption=Perintah pip yang umum]
# Upgrade pip ke versi terbaru
python -m pip install --upgrade pip

# Instal paket
python -m pip install requests

# Instal versi spesifik
python -m pip install requests==2.31.0

# Hapus paket
python -m pip uninstall requests

# Tampilkan paket yang terinstal
python -m pip list

# Informasi detail satu paket
python -m pip show requests

# Ekspor daftar dependensi ke requirements.txt
python -m pip freeze > requirements.txt

# Instal dari requirements.txt (untuk berbagi ke tim)
python -m pip install -r requirements.txt
\end{lstlisting}

\begin{konsep}
\textbf{Alur kerja standar proyek Python:}
\begin{enumerate}[leftmargin=*]
  \item Buat virtual environment: \texttt{python -m venv .venv}
  \item Aktifkan venv
  \item Instal dependensi: \texttt{pip install -r requirements.txt}
  \item Kerjakan proyek
  \item Jika menambah paket: \texttt{pip freeze > requirements.txt} (update)
\end{enumerate}
\end{konsep}

\section{Debugging dan Gaya Kode}

\begin{itemize}[leftmargin=*]
  \item Baca pesan error (\textit{traceback}) dari \textbf{bawah ke atas} --- baris paling bawah menunjuk ke sumber masalah.
  \item Tambahkan \texttt{print} diagnostik sementara atau gunakan debugger IDE\newline (VS Code, PyCharm).
  \item Ikuti \textbf{PEP 8}: indentasi 4 spasi, baris tidak lebih dari 79 karakter, nama konsisten (\texttt{snake\_case}).
  \item Gunakan \texttt{help(fungsi)} atau \texttt{dir(objek)} di REPL untuk eksplorasi cepat.
\end{itemize}

\begin{lstlisting}[caption=Teknik debugging sederhana dengan print]
# Debug dengan print (sementara)
def cari_maks(daftar):
    print(f"DEBUG: input={daftar}")   # hapus setelah debug
    if not daftar:
        return None
    maks = daftar[0]
    for x in daftar[1:]:
        if x > maks:
            maks = x
    print(f"DEBUG: maks={maks}")      # hapus setelah debug
    return maks
\end{lstlisting}

\section{Studi Kasus dan Contoh Kode}

Berikut adalah contoh-contoh program Python yang mengimplementasikan konsep-konsep pada modul ini.

\subsection{Kode: \texttt{utilitas.py}}
\begin{lstlisting}[language=Python,caption={\texttt{utilitas.py}}]
"""Modul utilitas contoh yang dapat diimpor atau dijalankan."""


def sapa_pengguna(nama: str) -> str:
    return f"Halo, {nama}! Selamat datang."

def main() -> None:
    # Blok ini digunakan untuk pengujian ketika dijalankan langsung
    print("Menguji modul utilitas...")
    print(sapa_pengguna("Budi"))

if __name__ == "__main__":
    main()
\end{lstlisting}

\subsection{Kode: \texttt{program\_utama.py}}
\begin{lstlisting}[language=Python,caption={\texttt{program\_utama.py}}]
"""Program utama yang mengimpor modul utilitas dan standar."""
import random
from utilitas import sapa_pengguna


def main() -> None:
    # Menggunakan fungsi dari modul buatan sendiri
    pesan = sapa_pengguna("Siti")
    print(pesan)
    
    # Menggunakan modul standar Python
    angka_acak = random.randint(1, 100)
    print(f"Angka keberuntungan Anda hari ini: {angka_acak}")

if __name__ == "__main__":
    main()
\end{lstlisting}

\section{Referensi sesi}

\begin{itemize}[leftmargin=*]
  \item \textbf{[8, 9]} \textit{The Python Tutorial} --- modul, \texttt{venv}, pustaka standar.
  \item \textbf{[14]} \textit{Dive Into Python 3} --- impor dan struktur paket.
\end{itemize}

\begin{aktivitas}
  \item Buat skrip \texttt{utilitas.py} dengan beberapa fungsi utilitas, tambahkan \texttt{if \_\_name\_\_ == "\_\_main\_\_"} dengan pengujian fungsi, lalu impor fungsi tersebut dari skrip lain.
  \item Gunakan \texttt{dir()} dan \texttt{help()} di REPL untuk menjelajahi modul \texttt{random} dan temukan tiga fungsi yang belum Anda ketahui.
  \item Buat virtual environment baru, instal satu paket ringan (\texttt{requests} jika jaringan lab mengizinkan, atau \texttt{colorama}), dan simpan \texttt{requirements.txt}.
  \item Sengaja buat bug \texttt{NameError} dan \texttt{TypeError}; perbaiki berdasarkan traceback.
  \item Format ulang potongan kode tidak PEP 8 menjadi lebih konsisten (nama variabel, spasi, panjang baris).
\end{aktivitas}

\begin{latihan}
  \item Jelaskan perbedaan \texttt{import mod} dan \texttt{from mod import fungsi}.
  \item Apa perbedaan antara \texttt{dir(math)} dan \texttt{help(math)}?
  \item Mengapa \texttt{if \_\_name\_\_ == "\_\_main\_\_"} penting saat membuat modul yang dapat diimpor?
  \item Apa manfaat \texttt{python -m pip} dibanding menjalankan \texttt{pip} langsung?
  \item Apa tujuan berkas \texttt{requirements.txt} dalam proyek tim?
  \item Refleksi: strategi debugging mana yang paling membantu Anda?
\end{latihan}

\begin{asesmen}
\textbf{Kuis}

\begin{enumerate}
  \item Perintah yang aman untuk menginstal paket ke interpreter yang sedang dipakai:
  \begin{enumerate}
    \item \texttt{python -m pip install nama\_paket}
    \item \texttt{pip install nama\_paket} (tanpa venv)
    \item \texttt{python install nama\_paket}
    \item \texttt{sudo pip install nama\_paket}
  \end{enumerate}
  \item Blok \texttt{if \_\_name\_\_ == "\_\_main\_\_"} \textbf{tidak} dieksekusi ketika:
  \begin{enumerate}
    \item Berkas dijalankan langsung dengan \texttt{python berkas.py}
    \item Berkas diimpor oleh modul lain
    \item Berkas dijalankan di Jupyter
    \item Berkas dijalankan di REPL dengan \texttt{exec()}
  \end{enumerate}
\end{enumerate}

\textbf{Proyek mini}: skrip \texttt{proyek\_mini.py} yang: (1) membaca konfigurasi JSON sederhana, (2) memvalidasi kunci yang diperlukan, (3) memproses data berdasarkan konfigurasi, (4) menulis log ke berkas. Gunakan \texttt{if \_\_name\_\_ == "\_\_main\_\_"} dan pisahkan fungsi-fungsi ke modul sendiri.
\end{asesmen}

\begin{checklist}
  \item Saya dapat mengimpor modul standar dengan benar (\texttt{import} dan \texttt{from ... import})
  \item Saya dapat menggunakan pola \texttt{if \_\_name\_\_ == "\_\_main\_\_"}
  \item Saya dapat menggunakan \texttt{dir()} untuk menjelajahi modul
  \item Saya dapat membuat dan mengaktifkan virtual environment
  \item Saya dapat menginstal paket, membuat dan menggunakan \texttt{requirements.txt}
  \item Saya dapat membaca traceback dan memperbaiki kesalahan sederhana
  \item Saya menerapkan minimal satu aturan gaya PEP 8 pada tugas
\end{checklist}

\begin{rangkuman}
Modul menyusun kode besar menjadi unit yang dapat diimpor. Pola \texttt{if \_\_name\_\_ == "\_\_main\_\_"} memungkinkan berkas berfungsi ganda sebagai skrip dan modul. \texttt{dir()} dan \texttt{help()} membantu eksplorasi API. Virtual environment mengisolasi dependensi antar proyek; \texttt{pip} mengelola dependensi pihak ketiga; \texttt{requirements.txt} memungkinkan reproduksi lingkungan.

\textbf{Referensi:} {\small Python Tutorial §6, §10, §12 \url{https://docs.python.org/3/tutorial/modules.html} dan \url{https://docs.python.org/3/tutorial/venv.html}}

\textbf{Kata kunci:} \code{import}, \code{\_\_name\_\_}, \code{dir()}, \code{help()}, \code{pip}, \code{venv}, \code{requirements.txt}, \code{PEP 8}
\end{rangkuman}

\end{document}
